%======================================================
\documentclass[12pt]{article}
%======================================================
\usepackage[utf8]{inputenc}
\usepackage[T1]{fontenc}
\usepackage[spanish,mexico]{babel}
\usepackage{lmodern}
\usepackage{amsmath,amssymb,amsthm,mathtools}
\usepackage{geometry,enumitem,booktabs,array,xcolor,graphicx,hyperref,fancyhdr,microtype,longtable,multicol,tikz}
%======================================================
\geometry{letterpaper,top=2.4cm,bottom=2.35cm,left=2.45cm,right=2.45cm}
\setlength{\parindent}{0pt}
\setlength{\parskip}{0.52em}
\setlength{\headheight}{15pt}
\setlength{\emergencystretch}{2em}
%======================================================
\definecolor{azul}{RGB}{22,72,120}
\hypersetup{colorlinks=true,linkcolor=azul,urlcolor=azul,pdftitle={Cálculo Diferencial - Unidad 4: Derivadas},pdfauthor={Carlos Ernesto Martinez}}
%======================================================
\pagestyle{fancy}
\fancyhf{}
\lhead{\small Tecnológico Nacional de México}
\rhead{\small Cálculo Diferencial}
\lfoot{\footnotesize Carlos Ernesto Martínez}
\cfoot{\footnotesize\thepage}
\rfoot{\footnotesize mrdgzcarlose@gmail.com}
\renewcommand{\headrulewidth}{0.4pt}
\renewcommand{\footrulewidth}{0.3pt}
\newtheorem{definicion}{Definición}[section]
\newtheorem{teorema}[definicion]{Teorema}
\newtheorem{proposicion}[definicion]{Proposición}
\newtheorem{corolario}[definicion]{Corolario}
\theoremstyle{remark}
\newtheorem{observacion}[definicion]{Observación}
\newtheorem{ejemplo}[definicion]{Ejemplo}
\newtheorem{Ejer}[definicion]{Ejercicio}
\newcommand{\R}{\mathbb{R}}
\newcommand{\C}{\mathbb{C}}
\newcommand{\vect}[1]{\boldsymbol{#1}}
\newcommand{\idea}[1]{\begin{center}\fcolorbox{azul}{white}{\parbox{0.91\linewidth}{#1}}\end{center}}
%======================================================
\title{\textbf{Notas para el curso de Cálculo Diferencial}\\[0.35cm]
\Large Unidad 4: Derivadas}
\author{Carlos Ernesto Martínez\\Tecnológico Nacional de México}
\date{Agosto--Diciembre 2026}
%======================================================
\begin{document}
\maketitle
\tableofcontents
\newpage
%=======================================================
\section*{Propósito y organización de la unidad}
%=======================================================
\textbf{Objetivo del temario.} Utilizar la definición de derivada para el análisis de funciones y el cálculo de derivadas.

\textbf{Organización.} La Unidad 4 comprende concepto y cálculo por definición; interpretación geométrica y física; propiedades y teoremas; reglas de derivación; funciones compuestas; y derivación implícita. El temario prevé 16 horas efectivas de clase y una hora para examen. El material desarrolla los conceptos antes de introducir los algoritmos, para que las reglas no aparezcan como fórmulas aisladas.

\section{De la razón de cambio promedio a la derivada}
Una función $y=f(x)$ describe cómo depende una cantidad de otra. Si el argumento varía de $a$ a $a+h$, la variación en la función es
\[
\Delta y=f(a+h)-f(a),\qquad \Delta x=h.
\]
Siempre que $h\ne0$, la \emph{razón de cambio promedio} de $f$ entre $a$ y $a+h$ es
\[
\frac{\Delta y}{\Delta x}=\frac{f(a+h)-f(a)}{h}.
\]
La derivada se obtiene haciendo que el intervalo de observación sea cada vez más pequeño, siempre que la razón anterior tenga límite finito.

\begin{definicion}[Derivada de una función en un punto]
Sea $f$ una función definida en un entorno de $a$. Su derivada en $a$ es
\[
\boxed{f'(a)=\lim_{h\to0}\frac{f(a+h)-f(a)}{h},}
\]
si el límite existe y es finito. En ese caso decimos que $f$ es \emph{derivable} o \emph{diferenciable} en $a$.
\end{definicion}

\begin{observacion}
La expresión $[f(a+h)-f(a)]/h$ se llama \emph{cociente incremental}. En su definición, $h\ne0$; escribir $h=0$ directamente produciría una división por cero. Debemos simplificar o transformar y luego calcular el límite.
\end{observacion}

\begin{ejemplo}[Derivada de $x^2$ en un punto]
Para $f(x)=x^2$ y un número real fijo $a$,
\begin{align*}
f'(a)&=\lim_{h\to0}\frac{(a+h)^2-a^2}{h}\\
&=\lim_{h\to0}\frac{2ah+h^2}{h}
=\lim_{h\to0}(2a+h)=\boxed{2a}.
\end{align*}
En particular, $f'(3)=6$. El cálculo no se basa en recordar la regla de la potencia: la hemos obtenido directamente de la definición.
\end{ejemplo}

\begin{definicion}[Función derivada]
Si la derivada existe en los puntos de un conjunto, definimos una nueva función $f'$ asignando a cada $x$ el valor
\[
\boxed{f'(x)=\lim_{h\to0}\frac{f(x+h)-f(x)}{h}.}
\]
También utilizamos $\displaystyle\frac{dy}{dx}$, $\displaystyle\frac{d}{dx}f(x)$ y $Df(x)$ para denotar la derivada respecto de $x$.
\end{definicion}

\begin{ejemplo}[Derivar un polinomio por definición]
Sea $f(x)=3x^2-2x+1$. Entonces
\begin{align*}
f(x+h)-f(x)
&=3(x+h)^2-2(x+h)+1-(3x^2-2x+1)\\
&=6xh+3h^2-2h.
\end{align*}
Dividiendo entre $h$ y tomando el límite,
\[
f'(x)=\lim_{h\to0}(6x+3h-2)=\boxed{6x-2}.
\]
\end{ejemplo}

\subsection*{Derivadas laterales y puntos no derivables}
La derivada derecha se calcula con $h\to0^+$, y la izquierda con $h\to0^-$. Para que exista $f'(a)$ como derivada bilateral ambas deben existir como números finitos y coincidir.

\begin{ejemplo}[Una esquina]
Para $f(x)=|x|$ en $a=0$,
\[
\frac{f(0+h)-f(0)}h=\frac{|h|}{h}
=\begin{cases}1,&h>0,\\-1,&h<0.\end{cases}
\]
La derivada derecha vale $1$, mientras que la izquierda vale $-1$. Por tanto, $\boxed{f'(0)\text{ no existe}}$. Sin embargo, $f$ sí es continua en $0$.
\end{ejemplo}

\begin{teorema}[Derivabilidad implica continuidad]
Si $f$ es derivable en $a$, entonces $f$ es continua en $a$.
\end{teorema}
\begin{proof}
Para $h\ne0$,
\[
f(a+h)-f(a)=h\,\frac{f(a+h)-f(a)}h.
\]
Como $f$ es derivable, el segundo factor tiende al número finito $f'(a)$ y el primero tiende a $0$. Por las leyes de límites, $f(a+h)-f(a)\to0$, es decir, $f(a+h)\to f(a)$. Esto es precisamente continuidad en $a$.
\end{proof}

\begin{observacion}
La implicación inversa es falsa: $|x|$ es continua en $0$ y no derivable allí. Una discontinuidad impide la derivabilidad, pero la continuidad no la garantiza. También puede faltar derivabilidad en cúspides, esquinas o puntos de tangente vertical.
\end{observacion}

\section{Interpretación geométrica de la derivada}
Para los puntos $P=(a,f(a))$ y $Q=(a+h,f(a+h))$ de la gráfica, la pendiente de la recta secante es
\[
m_{PQ}=\frac{f(a+h)-f(a)}h.
\]
Al acercar $Q$ a $P$, las pendientes de las secantes tienden, si existe el límite, a la pendiente de la tangente:
\[
\boxed{m_{\mathrm{tan}}=f'(a).}
\]

\begin{proposicion}[Recta tangente]
Si $f$ es derivable en $a$, la ecuación de su recta tangente en $(a,f(a))$ es
\[
\boxed{y-f(a)=f'(a)(x-a).}
\]
Cuando $f'(a)\ne0$, la recta normal, perpendicular a la tangente, tiene pendiente $-1/f'(a)$.
\end{proposicion}

\begin{ejemplo}[Tangente a una parábola]
Para $f(x)=x^2$ en $a=1$, conocemos $f(1)=1$ y $f'(1)=2$. Entonces
\[
y-1=2(x-1),\qquad\boxed{y=2x-1}.
\]
La tangente pasa por $(1,1)$ y tiene pendiente $2$. La normal tiene pendiente $-1/2$ y ecuación $y-1=-(x-1)/2$.
\end{ejemplo}

\begin{center}
\begin{tikzpicture}[scale=0.9,>=stealth]
\draw[->] (-0.3,0)--(3.6,0) node[right] {$x$};
\draw[->] (0,-1.1)--(0,4.5) node[above] {$y$};
\draw[thick,azul,domain=-.2:2.05,samples=80] plot (\x,{\x*\x}) node[above] {$y=x^2$};
\draw[thick,domain=.05:2.6] plot (\x,{2*\x-1});
\fill (1,1) circle (2.5pt) node[above left] {$(1,1)$};
\node[right] at (2.45,3.65) {$y=2x-1$};
\end{tikzpicture}
\end{center}

\subsection*{Significado de la pendiente}
Una derivada positiva indica que, localmente, pequeños incrementos positivos de $x$ producen incrementos de $f$ de primer orden positivos; una derivada negativa indica la variación contraria. Si $f'(a)=0$, la tangente es horizontal, pero ese hecho \emph{no} basta por sí solo para concluir que exista un máximo o un mínimo; esa aplicación se estudiará en la Unidad 5.

\section{Interpretación física de la derivada}
Si $s(t)$ expresa posición sobre una recta, su variación $s(t+h)-s(t)$ en el tiempo $h$ determina una velocidad media. El límite de estas velocidades medias es la velocidad instantánea:
\[
\boxed{v(t)=s'(t).}
\]
La aceleración instantánea es la razón de cambio de la velocidad:
\[
\boxed{a(t)=v'(t)=s''(t),}
\]
si ambas derivadas existen. Si $s$ está en metros y $t$ en segundos, $v$ se expresa en metros por segundo y $a$ en metros por segundo cuadrado.

\begin{ejemplo}[Posición, velocidad y aceleración]
Supongamos $s(t)=t^2+2t$ metros, para $t\ge0$. En un intervalo $[t,t+h]$, la velocidad media es
\begin{align*}
\frac{s(t+h)-s(t)}h
&=\frac{(t+h)^2+2(t+h)-(t^2+2t)}h\\
&=2t+h+2.
\end{align*}
Al tomar $h\to0$ obtenemos $v(t)=2t+2$ m/s. Derivando nuevamente, $a(t)=2$ m/s$^2$. En $t=3$ s, $v(3)=8$ m/s.
\end{ejemplo}

\begin{observacion}
La velocidad puede ser negativa cuando el movimiento ocurre en la dirección elegida como negativa. La rapidez es $|v(t)|$ y nunca es negativa. El desplazamiento es $s(b)-s(a)$; no se identifica automáticamente con la distancia total recorrida.
\end{observacion}

\section{Propiedades y reglas básicas de derivación}
Antes de memorizar reglas es conveniente justificar de dónde proceden. Las demostraciones siguientes se basan en el cociente incremental y las leyes de los límites vistas en la Unidad 3.

\begin{proposicion}[Constante, suma y múltiplo escalar]
Si $f$ y $g$ son derivables en $x$ y $c$ es constante, entonces
\[
\boxed{\frac{d}{dx}c=0,\qquad(f\pm g)'=f'\pm g',\qquad(cf)'=cf'.}
\]
\end{proposicion}
\begin{proof}
La función constante tiene incremento cero. Para la suma,
\[
\frac{(f+g)(x+h)-(f+g)(x)}h
=\frac{f(x+h)-f(x)}h+\frac{g(x+h)-g(x)}h.
\]
El límite de la derecha es $f'(x)+g'(x)$. Para el múltiplo escalar basta sacar $c$ fuera del cociente y luego del límite; la resta se obtiene con $c=-1$.
\end{proof}

\begin{teorema}[Regla del producto]
Si $f$ y $g$ son derivables,
\[\boxed{(fg)'=f'g+fg'.}\]
\end{teorema}
\begin{proof}
Sumamos y restamos el término $f(x)g(x+h)$ en el numerador:
\begin{align*}
&\frac{f(x+h)g(x+h)-f(x)g(x)}h\\
&\quad=\frac{f(x+h)-f(x)}h\,g(x+h)
+f(x)\frac{g(x+h)-g(x)}h.
\end{align*}
Al pasar al límite, $g(x+h)\to g(x)$ por continuidad, y los dos cocientes incrementales tienden a $f'(x)$ y $g'(x)$. Concluimos la fórmula.
\end{proof}

\begin{ejemplo}
Si $y=(x^2+1)(3x-2)$, aplicamos la regla del producto sin desarrollar previamente:
\[
y'=2x(3x-2)+(x^2+1)3.
\]
Simplificando, $\boxed{y'=9x^2-4x+3}$.
\end{ejemplo}

\begin{teorema}[Regla del cociente]
Si $f$ y $g$ son derivables y $g(x)\ne0$, entonces
\[
\boxed{\left(\frac f g\right)'=\frac{f'g-fg'}{g^2}.}
\]
\end{teorema}
\begin{proof}
Primero calculamos la derivada del recíproco:
\begin{align*}
\frac{\frac1{g(x+h)}-\frac1{g(x)}}h
&=\frac{g(x)-g(x+h)}{h\,g(x+h)g(x)}\\
&\longrightarrow-\frac{g'(x)}{[g(x)]^2}.
\end{align*}
Usando la regla del producto para $f\cdot(1/g)$ se obtiene la identidad indicada.
\end{proof}

\begin{ejemplo}
Para $y=(x^2+1)/(x-1)$, con $x\ne1$,
\begin{align*}
y'&=\frac{2x(x-1)-(x^2+1)}{(x-1)^2}\\
&=\boxed{\frac{x^2-2x-1}{(x-1)^2}}.
\end{align*}
\end{ejemplo}

\subsection*{Potencias y polinomios}
\begin{teorema}[Regla de la potencia, exponentes enteros positivos]
Para $n\in\mathbb{N}$, $n\ge1$,
\[
\boxed{\frac{d}{dx}x^n=nx^{n-1}.}
\]
\end{teorema}
\begin{proof}
El binomio de Newton permite escribir
\[
(x+h)^n-x^n=nx^{n-1}h+\binom n2x^{n-2}h^2+\cdots+h^n.
\]
Después de dividir entre $h$, todos los términos excepto $nx^{n-1}$ contienen una potencia positiva de $h$ y tienden a cero. Por tanto, el límite del cociente incremental es $nx^{n-1}$.
\end{proof}

\begin{corolario}
Si $p(x)=a_nx^n+\cdots+a_1x+a_0$, entonces
\[
p'(x)=na_nx^{n-1}+\cdots+2a_2x+a_1.
\]
La regla de la potencia se extiende a exponentes enteros negativos en $x\ne0$ y, mediante composición y las propiedades de potencias, a exponentes racionales en sus dominios de derivabilidad.
\end{corolario}

\begin{ejemplo}
\[
\frac{d}{dx}\left(\frac{3}{x^2}-2\sqrt x+5x\right)
=\frac{d}{dx}(3x^{-2}-2x^{1/2}+5x)
=\boxed{-6x^{-3}-x^{-1/2}+5},\quad x>0.
\]
\end{ejemplo}

\section{Derivadas de funciones elementales}
Las derivadas de las funciones elementales se obtienen combinando límites fundamentales con las reglas anteriores. Cada identidad debe aplicarse únicamente donde la expresión esté definida y sea derivable.

\subsection*{Funciones trigonométricas}
Recordemos que los argumentos trigonométricos se expresan en radianes al aplicar estas fórmulas.

\begin{ejemplo}[Derivada del seno a partir de límites]
Por la identidad $\sin(x+h)=\sin x\cos h+\cos x\sin h$,
\begin{align*}
\frac{\sin(x+h)-\sin x}{h}
&=\sin x\frac{\cos h-1}{h}+\cos x\frac{\sin h}{h}\\
&\longrightarrow 0\cdot\sin x+1\cdot\cos x.
\end{align*}
Así, $\boxed{(\sin x)'=\cos x}$.
\end{ejemplo}

\begin{proposicion}[Derivadas trigonométricas]
En los puntos donde las funciones están definidas,
\begin{align*}
(\sin x)'&=\cos x,&(\cos x)'&=-\sin x,\\
(\tan x)'&=\sec^2x,&(\cot x)'&=-\csc^2x,\\
(\sec x)'&=\sec x\tan x,&(\csc x)'&=-\csc x\cot x.
\end{align*}
\end{proposicion}

\begin{proof}[Deducción de la derivada de la tangente]
Como $\tan x=\sin x/\cos x$, y suponiendo $\cos x\ne0$,
\begin{align*}
(\tan x)'&=\frac{\cos x\cdot\cos x-\sin x(-\sin x)}{\cos^2x}\\
&=\frac{\cos^2x+\sin^2x}{\cos^2x}=\boxed{\sec^2x}.
\end{align*}
\end{proof}

\subsection*{Exponenciales y logaritmos}
\begin{proposicion}
Para $a>0$ con $a\ne1$ y para $x>0$ cuando se necesite,
\[
\boxed{(e^x)'=e^x,\quad(a^x)'=a^x\ln a,\quad(\ln x)'=\frac1x,\quad(\log_a x)'=\frac1{x\ln a}.}
\]
Las identidades pueden justificarse mediante los límites fundamentales de la exponencial y del logaritmo, o mediante la relación de función inversa.
\end{proposicion}

\begin{ejemplo}[Derivada del logaritmo por incremento]
Para $x>0$ y $h$ suficientemente pequeño, usando $\ln(x+h)-\ln x=\ln(1+h/x)$,
\begin{align*}
\frac{\ln(x+h)-\ln x}h
&=\frac1x\frac{\ln(1+h/x)}{h/x}\\
&\longrightarrow\frac1x,
\end{align*}
porque $\lim_{u\to0}\ln(1+u)/u=1$.
\end{ejemplo}

\begin{observacion}
Para $x<0$, la expresión $\ln x$ no es real. Sin embargo, $\ln|x|$ está definida para todo $x\ne0$ y satisface $\frac{d}{dx}\ln|x|=1/x$. Es importante distinguir entre dominio de definición y fórmula formal.
\end{observacion}

\section{Derivación de funciones compuestas: regla de la cadena}
Las funciones compuestas son frecuentes: primero calculamos $u=g(x)$ y, después, evaluamos $f(u)$. La regla de la cadena expresa cómo se combinan ambas razones de cambio.

\begin{teorema}[Regla de la cadena]
Si $g$ es derivable en $x$ y $f$ es derivable en $g(x)$, entonces $f\circ g$ es derivable en $x$ y
\[
\boxed{\frac{d}{dx}f(g(x))=f'(g(x))g'(x).}
\]
En notación $y=f(u)$, $u=g(x)$, también escribimos $\displaystyle\frac{dy}{dx}=\frac{dy}{du}\frac{du}{dx}$, entendiendo estas expresiones como derivadas relacionadas, no como una cancelación formal indiscriminada.
\end{teorema}

\begin{proof}[Justificación mediante incrementos]
La derivabilidad de $f$ en $u_0=g(x)$ significa que podemos escribir
\[
f(u_0+k)-f(u_0)=f'(u_0)k+r(k)k,
\]
donde $r(k)\to0$ cuando $k\to0$; definimos $r(0)=0$ para admitir también $k=0$. Tomemos $k=g(x+h)-g(x)$. Por continuidad de $g$, $k\to0$. Entonces
\begin{align*}
\frac{f(g(x+h))-f(g(x))}h
&=\bigl(f'(g(x))+r(k)\bigr)\frac{g(x+h)-g(x)}h.
\end{align*}
Al tomar $h\to0$ obtenemos $f'(g(x))g'(x)$. Esta demostración es válida incluso si algunos incrementos satisfacen $k=0$.
\end{proof}

\begin{ejemplo}[Potencia de un polinomio]
Para $y=(3x^2-1)^5$ identificamos la función interior $u=3x^2-1$ y la exterior $y=u^5$. Así,
\[
\frac{dy}{du}=5u^4,\qquad\frac{du}{dx}=6x,
\]
de modo que
\[
\boxed{y'=30x(3x^2-1)^4.}
\]
No es necesario desarrollar la quinta potencia.
\end{ejemplo}

\begin{ejemplo}[Composición trigonométrica]
Si $y=\sin(x^2+1)$, la función exterior es seno y la interior es $x^2+1$. Por tanto,
\[\boxed{y'=\cos(x^2+1)\,(2x).}\]
Si $y=\sin^3x$, interpretamos $(\sin x)^3$ y obtenemos $\boxed{y'=3\sin^2x\cos x}$.
\end{ejemplo}

\begin{ejemplo}[Composición logarítmica y exponencial]
En el intervalo donde $2x+1>0$,
\[
\frac{d}{dx}\ln(2x+1)=\frac{2}{2x+1}.
\]
Para todo real $x$,
\[
\frac{d}{dx}e^{x^2-3x}=e^{x^2-3x}(2x-3).
\]
El factor derivado de la función interior es indispensable.
\end{ejemplo}

\subsection*{Tabla de reglas para composiciones}
Si $u=u(x)$ es derivable y las expresiones están definidas,
\begin{align*}
(u^r)'&=ru^{r-1}u',&(\sqrt{u})'&=\frac{u'}{2\sqrt u},\\
(\sin u)'&=\cos u\,u',&(\cos u)'&=-\sin u\,u',\\
(\tan u)'&=\sec^2u\,u',&(e^u)'&=e^u u',\\
(a^u)'&=a^u\ln(a)\,u',&(\ln|u|)'&=\frac{u'}u\quad(u\ne0).
\end{align*}
Para exponentes reales $r$, la fórmula $u^r$ se aplica cuando $u>0$; otras extensiones requieren examinar su dominio.

\subsection*{Funciones trigonométricas inversas}
\begin{proposicion}
Para los valores de $x$ en los intervalos donde existen estas derivadas,
\[
\boxed{(\arcsin x)'=\frac1{\sqrt{1-x^2}},\quad
(\arccos x)'=-\frac1{\sqrt{1-x^2}},\quad
(\arctan x)'=\frac1{1+x^2}.}
\]
Las dos primeras fórmulas son válidas para $-1<x<1$, y la última para todo $x\in\R$.
\end{proposicion}

\begin{proof}[Deducción para $\arctan x$]
Sea $y=\arctan x$, de manera que $\tan y=x$ con $y\in(-\pi/2,\pi/2)$. Derivando la igualdad respecto de $x$ obtenemos $\sec^2y\,y'=1$. Como $\sec^2y=1+\tan^2y=1+x^2$,
\[\boxed{y'=\frac1{1+x^2}.}\]
\end{proof}

\section{Derivación implícita}
En una ecuación como $x^2+y^2=25$, la variable $y$ no aparece aislada. Sin embargo, cerca de ciertos puntos podemos describir $y$ como una función de $x$ y calcular su derivada sin despejarla explícitamente.

\begin{observacion}
Al derivar una expresión que contiene $y$, debemos tratar $y$ como $y(x)$. Por ejemplo,
\[
\frac{d}{dx}(y^2)=2y\,y',\qquad
\frac{d}{dx}(xy)=y+xy'.
\]
La primera igualdad usa la regla de la cadena; la segunda combina producto y cadena.
\end{observacion}

\begin{ejemplo}[Circunferencia]
Para $x^2+y^2=25$, derivamos a ambos lados:
\[
2x+2y\,y'=0.
\]
Despejamos la derivada:
\[
\boxed{\frac{dy}{dx}=-\frac{x}{y},\qquad y\ne0.}
\]
En el punto $(3,4)$ la pendiente es $-3/4$ y la recta tangente es
\[
y-4=-\frac34(x-3).
\]
En $(5,0)$ no podemos utilizar esta fórmula, porque la circunferencia tiene allí tangente vertical y no define localmente a $y$ como función derivable de $x$.
\end{ejemplo}

\begin{ejemplo}[Producto y potencia en una relación implícita]
Dada la relación $xy+y^2=6$, derivamos:
\[
\underbrace{y+xy'}_{(xy)'}+\underbrace{2yy'}_{(y^2)'}=0.
\]
Agrupando los términos que contienen $y'$,
\[
(x+2y)y'=-y,
\]
de donde
\[
\boxed{y'=-\frac{y}{x+2y}},\qquad x+2y\ne0.
\]
\end{ejemplo}

\begin{ejemplo}[Exponencial en derivación implícita]
Sea $e^y+xy=2$. Entonces
\[
e^y y'+y+xy'=0.
\]
Factorizamos $y'$ y despejamos:
\[
(e^y+x)y'=-y,\qquad
\boxed{y'=-\frac{y}{e^y+x}},\quad e^y+x\ne0.
\]
\end{ejemplo}

\begin{proposicion}[Fórmula implícita y sus condiciones]
Si $F(x,y)=0$ describe localmente una función derivable $y=y(x)$ y $F$ posee derivadas parciales continuas cerca del punto, con $F_y\ne0$, entonces
\[
\boxed{y'=-\frac{F_x}{F_y}.}
\]
La condición $F_y\ne0$ es una condición suficiente local para poder despejar diferenciablemente $y$ en función de $x$; si falla, la fórmula no se puede aplicar directamente.
\end{proposicion}

\begin{observacion}
Esta fórmula resume el procedimiento anterior, pero en este curso conviene practicar primero la derivación término a término. El método muestra de manera explícita dónde se utilizan las reglas del producto y de la cadena.
\end{observacion}

%=======================================================
\section{Ejemplos complementarios y comprobaciones}
%=======================================================
\subsection*{Dos formas equivalentes de escribir la derivada}
Cuando $x$ se aproxima a $a$, el incremento $h=x-a$ tiende a cero. Por sustitución de variables, la definición también puede escribirse
\[
\boxed{f'(a)=\lim_{x\to a}\frac{f(x)-f(a)}{x-a}},
\]
si existe el límite finito. La segunda forma puede resultar más cómoda cuando ya conocemos $f(x)$ y queremos estudiar un punto concreto; la primera es especialmente útil para desarrollar expresiones como $f(a+h)$.

\begin{ejemplo}[Derivada de la raíz cuadrada por definición]
Para $f(x)=\sqrt{x}$ y $a>0$,
\begin{align*}
f'(a)&=\lim_{h\to0}\frac{\sqrt{a+h}-\sqrt a}{h}\\
&=\lim_{h\to0}\frac{(\sqrt{a+h}-\sqrt a)(\sqrt{a+h}+\sqrt a)}{h(\sqrt{a+h}+\sqrt a)}\\
&=\lim_{h\to0}\frac{1}{\sqrt{a+h}+\sqrt a}
=\boxed{\frac1{2\sqrt a}}.
\end{align*}
La racionalización transforma la indeterminación $0/0$ en un cociente cuyo límite se obtiene directamente. Se trata de la misma técnica que estudiamos en la Unidad 3.
\end{ejemplo}

\subsection*{Continuidad, esquinas y tangentes verticales}
\begin{ejemplo}
Consideremos $f(x)=x^{1/3}$, continua en todo $\R$. En $x=0$ el cociente incremental es
\[
\frac{f(h)-f(0)}h=\frac{h^{1/3}}h=\frac1{|h|^{2/3}},\qquad h\ne0.
\]
Al tomar $h\to0$ el cociente crece hacia $+\infty$; no existe una derivada finita en el origen. Geométricamente, la curva presenta una tangente vertical. Este comportamiento es distinto de la esquina de $|x|$, donde las pendientes laterales son finitas pero diferentes.
\end{ejemplo}

\begin{ejemplo}
Para $g(x)=x^{2/3}$ tenemos
\[
\frac{g(h)-g(0)}h=\frac{|h|^{2/3}}h=
\begin{cases}|h|^{-1/3},&h>0,\\-|h|^{-1/3},&h<0.\end{cases}
\]
La pendiente del lado derecho tiende a $+\infty$, mientras que la del lado izquierdo tiende a $-\infty$. Hay una cúspide y tampoco existe derivada finita. Aunque ambos ejemplos son continuos, ninguno es derivable en el origen.
\end{ejemplo}

\subsection*{Derivada del coseno a partir de identidades trigonométricas}
Por la identidad $\cos(x+h)=\cos x\cos h-\sin x\sin h$,
\begin{align*}
\frac{\cos(x+h)-\cos x}{h}
&=\cos x\frac{\cos h-1}{h}-\sin x\frac{\sin h}{h}\\
&\longrightarrow \cos x\cdot0-\sin x\cdot1=-\sin x.
\end{align*}
De esta forma, $\boxed{(\cos x)'=-\sin x}$ se deduce de los mismos dos límites trigonométricos empleados para derivar el seno.

\subsection*{Deducción de la derivada de $\arcsin x$}
Sea $y=\arcsin x$, con $-1<x<1$ y $y\in[-\pi/2,\pi/2]$. La identidad $\sin y=x$ implica, después de derivar,
\[
\cos y\,y'=1,\qquad y'=\frac1{\cos y}.
\]
En el intervalo principal del arcoseno, $\cos y\ge0$. Como $\cos^2y=1-\sin^2y=1-x^2$, resulta
\[
\boxed{\frac{d}{dx}\arcsin x=\frac1{\sqrt{1-x^2}}},\qquad -1<x<1.
\]
En los extremos $x=\pm1$ el denominador se anula y no existe una derivada bilateral finita.

\subsection*{Elección de una técnica sin operaciones innecesarias}
\begin{ejemplo}
Para $y=(x^2+1)^3$, desarrollar la potencia produciría un polinomio de grado seis. La regla de la cadena permite conservar la estructura:
\[
y'=3(x^2+1)^2(2x)=\boxed{6x(x^2+1)^2}.
\]
En cambio, si $y=(x+1)(x-1)$, desarrollar primero produce $x^2-1$, cuya derivada es $2x$. En este caso ambas técnicas son breves. La elección adecuada depende de la expresión y del objetivo del ejercicio.
\end{ejemplo}

\subsection*{Derivación implícita y verificación mediante despeje}
\begin{ejemplo}
Dada la curva $xy=6$ y suponiendo $x\ne0$, podemos despejar $y=6/x$, cuya derivada es $y'=-6/x^2$. Por derivación implícita, en cambio,
\[
xy'= -y,\qquad y'=-\frac{y}{x}.
\]
Al sustituir $y=6/x$ recuperamos $y'=-6/x^2$. Las dos expresiones coinciden. En el punto $(2,3)$ la pendiente es $-3/2$ y la recta tangente es
\[
\boxed{y-3=-\frac32(x-2).}
\]
La comprobación es útil para fijar la regla del producto y el tratamiento de $y$ como función de $x$.
\end{ejemplo}

\section{Síntesis de métodos y comprobaciones}
Para derivar una expresión es aconsejable: identificar primero su dominio; distinguir si es suma, producto, cociente o composición; aplicar las reglas sólo después de reconocer esa estructura; realizar simplificaciones razonables al final; y comprobar un resultado sencillo mediante la definición, un desarrollo algebraico o una derivación alternativa cuando sea posible. En derivación implícita es indispensable colocar el factor $y'$ al derivar cada término que depende de $y$.

%=======================================================
\section{Banco de ejercicios de la Unidad 4}
%=======================================================
Los reactivos se organizan de manera gradual. Los primeros bloques requieren explícitamente la definición; los siguientes practican reglas, composiciones y derivación implícita. Se recomienda presentar el procedimiento y no sólo la expresión final. La numeración permanece continua entre bloques sin agregar prefijos a los números de los reactivos.

\begin{Ejer}\textbf{Concepto de derivada y cálculo mediante la definición}
\begin{enumerate}[series=banco]
\item Utilice la definición para calcular $f'(x)$ si $f(x)=3x+2$.
\item Utilice la definición para calcular $f'(x)$ si $f(x)=x^2$.
\item Calcule $f'(2)$ mediante el cociente incremental si $f(x)=x^2+1$.
\item Mediante la definición, derive $f(x)=2x^2-x$.
\item Mediante la definición, derive $f(x)=x^3$.
\item Mediante la definición, calcule $f'(1)$ si $f(x)=1/x$, $x\ne0$.
\item Utilice la definición para obtener $f'(x)$ si $f(x)=1/x$, $x\ne0$.
\item Utilice la definición para calcular $f'(x)$ si $f(x)=\sqrt{x}$, $x>0$.
\item Determine la derivada de $f(x)=|x|$ en $x=0$ mediante límites laterales.
\item Determine si $f(x)=\begin{cases}x^2,&x\le1,\\2x-1,&x>1\end{cases}$ es derivable en $1$.
\end{enumerate}
\end{Ejer}

\begin{Ejer}\textbf{Tangentes, normales y razones de cambio}
\begin{enumerate}[resume=banco]
\item Encuentre la recta tangente a $y=x^2$ en $x=2$.
\item Encuentre la recta tangente a $y=3x-1$ en $x=1$.
\item Encuentre la recta tangente a $y=x^3$ en $x=-1$.
\item Encuentre la recta normal a $y=x^2+1$ en el punto $(1,2)$.
\item Determine la pendiente de la secante de $y=x^2$ entre $x=1$ y $x=1+h$ y calcule su límite.
\item Si $s(t)=t^2+3t$ metros, obtenga la velocidad media en $[1,1+h]$ y la instantánea en $t=1$.
\item Para $s(t)=2t^2-t+1$, encuentre $v(t)$ y $a(t)$.
\item Si $s(t)=t^3-3t$ para $t\ge0$, calcule la velocidad instantánea en $t=2$.
\item Explique en qué se distinguen velocidad y rapidez cuando $v(t)<0$.
\item Determine la recta tangente a $f(x)=1/x$ en $x=1$ y verifique su pendiente a partir de la definición.
\end{enumerate}
\end{Ejer}

\begin{Ejer}\textbf{Reglas básicas, productos y cocientes}
\begin{enumerate}[resume=banco]
\item Derive $f(x)=5x^4-3x^2+7$.
\item Derive $f(x)=2x^5-x^3+4x-1$.
\item Derive $f(x)=3/x^2+2/x$ e indique el dominio.
\item Derive $f(x)=\sqrt x+x^3$, $x>0$.
\item Derive $f(x)=(x+1)(x^2-2)$ utilizando la regla del producto.
\item Derive $f(x)=(2x-1)(x^3+1)$ utilizando la regla del producto.
\item Derive $f(x)=(x^2+1)/(x+1)$ e indique el dominio.
\item Derive $f(x)=(3x-2)/(x^2+1)$.
\item Derive $f(x)=x^2/(x-2)$.
\item Derive $f(x)=(x^2-1)/(x^2+1)$.
\item Compruebe mediante desarrollo algebraico que la derivada de $(x+2)(x-3)$ coincide con la obtenida usando el producto.
\item Justifique por qué no es correcta, en general, la identidad $(fg)'=f'g'$; proporcione un contraejemplo.
\end{enumerate}
\end{Ejer}

\begin{Ejer}\textbf{Derivadas trigonométricas, exponenciales y logarítmicas}
\begin{enumerate}[resume=banco]
\item Derive $f(x)=3\sin x-2\cos x$.
\item Derive $f(x)=\tan x+\sec x$.
\item Derive $f(x)=x\sin x$.
\item Derive $f(x)=\sin x/\cos x$ usando la regla del cociente.
\item Derive $f(x)=e^x+3^x$.
\item Derive $f(x)=xe^x$.
\item Derive $f(x)=\ln x+x^2$, $x>0$.
\item Derive $f(x)=(\ln x)/x$, $x>0$.
\item Derive $f(x)=\arctan x+x$, indicando el dominio de la derivada.
\item Derive $f(x)=\arcsin x$ e indique en qué puntos existe una derivada finita.
\item Derive $f(x)=\sec x\tan x$ donde ambas funciones estén definidas.
\item Derive $f(x)=\ln|x|$ por separado en $x>0$ y $x<0$.
\end{enumerate}
\end{Ejer}

\begin{Ejer}\textbf{Regla de la cadena: composición de funciones}
\begin{enumerate}[resume=banco]
\item Derive $f(x)=(2x+1)^5$.
\item Derive $f(x)=(x^2+3)^4$.
\item Derive $f(x)=\sqrt{x^2+1}$.
\item Derive $f(x)=\sin(3x)$.
\item Derive $f(x)=\cos(x^2)$.
\item Derive $f(x)=\tan(2x-1)$.
\item Derive $f(x)=e^{x^2+1}$.
\item Derive $f(x)=\ln(1+x^2)$.
\item Derive $f(x)=\sin^3x$.
\item Derive $f(x)=(\sin x+1)^4$.
\item Derive $f(x)=\ln|x^2-4|$ en los puntos donde está definida.
\item Derive $f(x)=\sqrt{1+\sin x}$ en los puntos donde la raíz sea positiva.
\item Derive $f(x)=e^{\sin(x^2)}$.
\item Derive $f(x)=\arctan(2x+1)$.
\end{enumerate}
\end{Ejer}

\begin{Ejer}\textbf{Productos, cocientes y composiciones combinadas}
\begin{enumerate}[resume=banco]
\item Derive $f(x)=x^2\sin(3x)$.
\item Derive $f(x)=(x+1)e^{2x}$.
\item Derive $f(x)=(x^2+1)\cos(x^2)$.
\item Derive $f(x)=\ln(1+x)/x$, $x>0$.
\item Derive $f(x)=e^x/(x^2+1)$.
\item Derive $f(x)=\sqrt{x^2+1}/(x+1)$, $x\ne-1$.
\item Derive $f(x)=x\ln(x^2+1)$.
\item Derive $f(x)=\sin x/(1+\cos x)$ en los puntos de su dominio.
\item Derive $f(x)=(x^2-1)^3/(x+2)$.
\item Derive $f(x)=e^{2x}\cos(3x)$.
\item Derive $f(x)=(\ln x)^2$, $x>0$.
\item Derive $f(x)=\sin(\sqrt{x})$, $x>0$.
\end{enumerate}
\end{Ejer}

\begin{Ejer}\textbf{Derivación implícita y rectas tangentes}
\begin{enumerate}[resume=banco]
\item Derive implícitamente $x^2+y^2=9$ y despeje $y'$ donde sea posible.
\item Derive implícitamente $xy=4$ y obtenga $y'$.
\item Derive implícitamente $x^2+xy+y^2=7$.
\item Derive implícitamente $xy+y^2=3$.
\item Derive implícitamente $y^3+x^2y=4$.
\item Derive implícitamente $\sin y+x=1$.
\item Derive implícitamente $e^y+x^2=3$.
\item Derive implícitamente $\ln y+xy=2$, suponiendo $y>0$.
\item Para $x^2+y^2=25$, encuentre la recta tangente en $(3,4)$.
\item Para $xy=6$, encuentre la recta tangente en $(2,3)$.
\item Derive implícitamente $\sin(x+y)=x$.
\item Derive implícitamente $e^{xy}+y=2$.
\end{enumerate}
\end{Ejer}

\begin{Ejer}\textbf{Integradores: interpretación y elección de técnica}
\begin{enumerate}[resume=banco]
\item Decida si $f(x)=|x-2|$ es continua y derivable en $x=2$; justifique mediante límites.
\item Encuentre $f'(1)$ por definición si $f(x)=x^2+2x$ y compare con las reglas.
\item Calcule la derivada de $y=(x^2+1)^3\sin x$ indicando dónde usa producto y dónde cadena.
\item Determine la recta tangente a $y=\ln x$ en $x=1$.
\item Si $s(t)=t^3-6t^2+9t$ metros, determine $v(1)$ y $a(1)$ con unidades.
\item Para $x^2+xy+y^2=12$, obtenga la pendiente de la curva en $(2,2)$.
\item Compruebe que las dos expresiones de $d(\tan x)/dx$ obtenidas por cociente y por regla directa coinciden.
\item Explique por qué la fórmula $(f\circ g)'=f'\circ g$ es, en general, falsa; utilice $f(u)=u^2$, $g(x)=3x$.
\item Derive $f(x)=\ln(\sqrt{1+x^2})$ utilizando primero logaritmos y después la cadena; compare.
\item Investigue a partir de la definición por qué $f(x)=x^{2/3}$ no es derivable en $0$ aunque es continua allí.
\end{enumerate}
\end{Ejer}

\end{document}
